\documentclass[11pt]{article}
\usepackage[a4paper,margin=2.5cm]{geometry}
\usepackage{amsmath,amssymb,amsthm}
\usepackage{hyperref}

\newtheorem{theorem}{Theorem}
\newtheorem{conjecture}[theorem]{Conjecture}
\newtheorem{lemma}[theorem]{Lemma}
\newtheorem{remark}[theorem]{Remark}

\title{Disproof of TLMC Conjecture 00000000122\\
\large Prime values of $q$-Catalan numbers}
\author{TLMC submission}
\date{September 2026}

\begin{document}
\maketitle

\begin{conjecture}[TLMC 00000000122]
For every $n \ge 2$ there exist infinitely many primes $q$ such that the
$n$-th $q$-Catalan number $C_n(q)$ is prime.
\end{conjecture}

\begin{theorem}[Verdict]\label{thm:verdict}
Conjecture~\ref{thm:verdict} is \textbf{false}. The instance $n = 2$ admits at
most one prime value of $q$ under both standard readings of the
$q$-Catalan numbers.
\end{theorem}

\section{The attack: $n = 2$}

\subsection{Standard (Gaussian-binomial) reading}

Under the standard definition
\[
C_n(q) \;=\; \frac{\begin{bmatrix} 2n \\ n \end{bmatrix}_q}{[n+1]_q},
\qquad
\begin{bmatrix} a \\ b \end{bmatrix}_q = \prod_{i=1}^{b} \frac{1-q^{a-b+i}}{1-q^i},
\]
one computes
\[
\begin{bmatrix} 4 \\ 2 \end{bmatrix}_q
= \frac{(1-q^3)(1-q^4)}{(1-q)(1-q^2)}
= \frac{[3]_q\,[4]_q}{[1]_q\,[2]_q}
= (1+q^2)(1+q+q^2),
\]
so that, since $[3]_q = 1+q+q^2$,
\begin{equation}\label{eq:C2}
C_2(q) \;=\; \frac{(1+q^2)(1+q+q^2)}{1+q+q^2} \;=\; q^2 + 1 .
\end{equation}

\begin{lemma}\label{lem:standard}
If $q$ is an odd integer with $q \ge 3$, then $q^2 + 1$ is not prime.
\end{lemma}

\begin{proof}
$q$ odd implies $q^2$ odd, hence $q^2 + 1$ even, i.e.\ $2 \mid q^2+1$;
concretely $q = 2j+1$ gives $q^2+1 = 2\,(2j^2+2j+1)$. Moreover $q \ge 3$
gives $q^2 + 1 \ge 10 > 2$. A number divisible by $2$ and strictly larger
than $2$ is not prime.
\end{proof}

Every odd prime $q$ satisfies the hypotheses of
Lemma~\ref{lem:standard}. Hence the only prime $q$ for which $C_2(q)$ can be
prime is $q = 2$; and indeed $C_2(2) = 5$ \emph{is} prime. The $n = 2$
instance therefore has exactly one witness --- not infinitely many.

\subsection{Carlitz reading}

Under the Carlitz recurrence $C_0(q) = 1$,
$C_{n+1}(q) = \sum_{k=0}^{n} q^k C_k(q) C_{n-k}(q)$, one gets
$C_1(q) = 1$ and
\[
C_2(q) = C_0 C_1 + q\, C_1 C_0 = 1 + q .
\]
The identical argument applies: for odd prime $q$, $q + 1$ is even and
$\ge 4$, hence composite, and $C_2(2) = 3$ is prime. Exactly one witness.

\begin{theorem}[Exact classification for $n = 2$]\label{thm:classification}
Let $q$ be prime. Under the standard reading, $C_2(q)$ is prime if and only
if $q = 2$ (with $C_2(2) = 5$). Under the Carlitz reading, $C_2(q)$ is prime
if and only if $q = 2$ (with $C_2(2) = 3$).
\end{theorem}

\section{Boundary of the disproof}

\begin{remark}
The counterexample is the single instance $n = 2$. Nothing here excludes the
conjecture for any fixed $n \ge 3$: the falsity is a failure of the
universal quantifier over $n$, not of the infinitude phenomenon for each
individual degree. The surviving witness $q = 2$ is genuine in both
readings, so the failure is exactly ``one prime witness instead of
infinitely many''.
\end{remark}

\section{Lean certificate}

The attack is formalized in core Lean~4 (no Mathlib), toolchain
\texttt{leanprover/lean4:v4.33.1}, in the file \texttt{lean4/Main.lean}.
Headline statements:
\begin{itemize}
  \item \texttt{tlmc122\_disproof\_standard}: for every prime $q$,
        $\mathrm{IsPrime}(C_2(q)) \to q = 2$ with $C_2(q) = q*q + 1$;
  \item \texttt{tlmc122\_disproof\_carlitz}: the same for $C_2(q) = q + 1$;
  \item \texttt{C2\_prime\_iff\_eq\_two}, \texttt{C2Carlitz\_prime\_iff\_eq\_two}:
        the exact classification of Theorem~\ref{thm:classification},
        including the primality of $5$ and $3$.
\end{itemize}
The development uses a self-contained primality predicate and a
decomposition-based parity API, deliberately avoiding \texttt{omega},
\texttt{simp} and the core mod/div lemmas (which carry \texttt{propext} or
\texttt{Quot.sound} in this toolchain). The audit \texttt{lean4/Check.lean}
runs \texttt{\#print axioms} on all $21$ theorems; every one reports
\emph{does not depend on any axioms}. There is no \texttt{sorry}.

\section{Reproduction}

\texttt{reproduce.py} recomputes everything independently with exact integer
polynomial arithmetic:
\begin{enumerate}
  \item the Gaussian binomial $\begin{bmatrix} 4 \\ 2 \end{bmatrix}_q$ by the
        product formula with exact division, confirming $C_2(q) = q^2+1$;
  \item the Carlitz polynomials through $C_3$, confirming
        $C_2(q) = 1+q$;
  \item an enumeration of all primes $q < 500$: in both readings the only
        prime $q$ with $C_2(q)$ prime is $q = 2$.
\end{enumerate}
The script asserts each of these facts and exits non-zero on any mismatch.

\end{document}
